% Generated draft: requires mathematical and source review.
\documentclass[11pt]{article}
\usepackage[margin=25mm]{geometry}
\usepackage{fontspec}
\setmainfont{Latin Modern Roman}
\usepackage{amsmath,amssymb,mathtools,unicode-math}
\setmathfont{Latin Modern Math}
\usepackage{enumitem,needspace,xurl,hyperref}
\setlength{\parindent}{0pt}
\setlength{\parskip}{5pt}
\setlength{\emergencystretch}{3em}
\newcommand{\N}{\mathbb{N}}
\newcommand{\Z}{\mathbb{Z}}
\newcommand{\Q}{\mathbb{Q}}
\newcommand{\R}{\mathbb{R}}
\newcommand{\C}{\mathbb{C}}
\newcommand{\F}{\mathbb{F}}
\newcommand{\A}{\mathbb{A}}
\newcommand{\PP}{\mathbb{P}}
\newcommand{\E}{\mathbb{E}}
\newcommand{\eps}{\varepsilon}
\newcommand{\dd}{\,\mathrm{d}}
\begin{document}
\section*{UnsolvedMath ID: 30001214}
\textbf{Source problem code:} catalogue entry
\subsection{Definitions and mathematical statement}
Use the usual commutative-algebra interpretation of a complex affine Poisson algebra: a commutative unital finitely generated \(\C\)-algebra \(R\) equipped with a \(\C\)-bilinear Lie bracket \(\{-,-\}\) satisfying
\[
\{a,bc\}=\{a,b\}c+b\{a,c\}\qquad(a,b,c\in R).
\]
A Poisson ideal is an ideal \(J\subseteq R\) such that \(\{R,J\}\subseteq J\). For \(\mathfrak m\in\operatorname{MaxSpec}R\), its Poisson core is the largest Poisson ideal contained in \(\mathfrak m\), namely
\[
\mathcal P(\mathfrak m)=\sum_{\substack{J\subseteq\mathfrak m\\J\text{ a Poisson ideal}}}J.
\]
Equip \(X=\operatorname{MaxSpec}R\) with the Zariski topology, whose closed sets are
\[
V_X(I)=\{\mathfrak n\in X:I\subseteq\mathfrak n\}\qquad(I\subseteq R\text{ an ideal}).
\]
Define \(\mathfrak m\sim\mathfrak n\) if \(\mathcal P(\mathfrak m)=\mathcal P(\mathfrak n)\). The symplectic core through \(\mathfrak m\) is the equivalence class
\[
\mathcal C(\mathfrak m)=\{\mathfrak n\in X:\mathcal P(\mathfrak n)=\mathcal P(\mathfrak m)\}.
\]
\textbf{Question.} For every complex affine Poisson algebra \(R\) and every \(\mathfrak m\in X\), is \(\mathcal C(\mathfrak m)\) locally closed in \(X\)? Explicitly, must there exist an ideal \(I\subseteq R\) and a Zariski-open subset \(U\subseteq X\) such that
\[
\mathcal C(\mathfrak m)=U\cap V_X(I)?
\]
The supplied record does not specify whether \(R\) must be reduced or an integral domain. Neither restriction is imposed in the displayed formulation; whether either was intended remains a source qualification requiring review.
\subsection{Short English statement}
For every complex affine Poisson algebra, group maximal ideals according to their largest contained Poisson ideal. Must each resulting symplectic core be locally closed in the Zariski topology? The source does not specify reducedness or irreducibility assumptions.
\subsection{Sources}
\begin{itemize}
\item Supplied local UnsolvedMath registry record, permanent ID \(30001214\), title \emph{Local Closedness of Symplectic Cores}, exact source problem code \texttt{catalogue entry}, in the Oberwolfach Reports Open Problems collection. The supplied proposed year is 2009; no individual proposer is identified.
\item Supplied report citation: \emph{Mini-Workshop: Non-Negativity is a Quantum Phenomenon} (2009), Oberwolfach Reports, pp.~767--800, DOI \texttt{10.4171/owr/2009/14}. \url{https://doi.org/10.4171/owr/2009/14}. This is also the supplied literature source labeled ``original OWR report.''
\item Supplied literature source labeled ``Brown--Gordon primary paper'': \url{https://arxiv.org/abs/math/0201042}.
\item Supplied literature source labeled ``Poisson Dixmier--Moeglin criteria'': \url{https://arxiv.org/abs/1908.06542}.
\item Literature assessment dated exactly \texttt{2026-08-21}: the supplied curation labels the question open and states that general results provide criteria and special cases but do not establish local closedness for all Poisson varieties. This is supplied dated curation, not independently verified here. The record expressly requests independent verification of time-sensitive status and cited claims; its reference to all Poisson varieties does not itself specify an extension of the affine question.
\end{itemize}
\end{document}
